% sections/discussion.tex -- the conclusion (paper_v2_vn). Since 20 November (N26) there is no
% Discussion or Threats section: their main points are in Section 7 (why the searched sets lose, why
% GAT leads) and in the conclusion below (the cheapest choice, who should judge a search, the
% limitations). The two minor limitations, in an appendix from N26, are back here since N31. Every
% number is \val{} from numbers.tex.
\section{Kết luận}
\label{sec:conclusion}

Trên một phiên bản Amazon-Meta đã làm sạch với ba kiểu nút và \val{search-cands} meta-path ứng
viên, kết quả của tìm kiếm meta-path tự động khá gần với thiết kế thủ công cho gợi ý sản phẩm dựa
trên HAN: các tập mà một phép tìm kiếm gồm \val{search-runs-total} lần chạy chọn ra kém các tập thủ
công \val{han-6-hit-relabs}--\val{han-4-hit-relabs}\% theo độ đo đã cố định từ trước và ngang bằng
chúng khi xếp hạng toàn bộ. Sự thua kém chủ yếu nằm ở những sản phẩm mà các meta-path danh mục của
các tập tìm được không đến được. Vì vậy, nên chọn cách nào tùy thuộc vào bối cảnh. Nếu có chuyên
gia để thiết kế được meta-path trong thời gian ngắn hơn thời gian tìm kiếm, ở đây sẽ khoảng
\val{search-total-h}~giờ~GPU, thiết kế thủ công là lựa chọn tốt hơn. Nếu không có chuyên gia,
tìm kiếm tự động cho các tập gần với thiết kế thủ công mà không cần kiến thức chuyên môn. Tuy
nhiên, đối chứng ngẫu nhiên cho thấy trên lược đồ nhỏ này việc chọn meta-path ít ảnh hưởng: các
tập chọn ngẫu nhiên từ cùng các ứng viên vốn không tốn chi phí, đạt kết quả ít nhất ngang bằng
cả hai cách (Mục~\ref{sec:res-random}). Lợi thế của tìm kiếm (nếu có) nhiều khả năng nằm ở những
lược đồ lớn hơn, nơi cả thiết kế thủ công lẫn chọn ngẫu nhiên đều khó, điều mà nghiên cứu này chưa
kiểm chứng. Một GAT không phân biệt kiểu vượt cả hai mô hình dùng meta-path.

Bài học chính là một phương pháp tìm kiếm cho một mô hình gợi ý nên được đánh giá bằng chính mô hình
gợi ý đó, như RMS đã làm \citep{ning2022automatic}, hoặc ít nhất được kiểm tra so với các tập chọn
ngẫu nhiên từ chính các ứng viên của nó. Các tập mà mạng riêng của LMSPS đánh giá là tốt nhất không
phải là các tập tốt nhất đối với HAN. Mức cải thiện của LMSPS trên HAN
\citep[Bảng~4]{li2024longrange} được đo trên bài toán phân loại nút, nơi mỗi nút được mã hóa từ vùng
lân cận của nó. Trong bài toán dự đoán liên kết của chúng tôi, phần lớn tham số của HAN là các
embedding người dùng tự do mà không meta-path nào định hình.

\paragraph{Hạn chế.} Nghiên cứu dùng một bộ dữ liệu, một bài toán và một lược đồ nhỏ, chỉ với các
meta-path phía sản phẩm. Phép chia ngẫu nhiên duy nhất bỏ qua ngày nhận xét, điều cho phép mô
hình học từ các tương tác về sau và có thể làm thay đổi thứ hạng của các mô hình
\citep{ji2023leakage,meng2020splitting}. Mọi tập dùng chung một cấu hình không tinh chỉnh. Khi giai
đoạn khởi động trọng số meta-path ngắn hơn, tập tìm được lại vượt tập thủ công ở sáu meta-path,
trong một phép kiểm tra chỉ dùng một seed (Mục~\ref{sec:res-sensitivity}). Phép tìm kiếm của chúng
tôi nhỏ hơn so với bài báo gốc vì được định cỡ cho một GPU. Các lần chạy của nó đồng thuận hơn mức
ngẫu nhiên, nhưng một phép tìm kiếm lớn hơn có thể chọn khác. Các tập thủ công được thiết kế theo các nguyên
tắc nêu rõ chứ không phải bởi một chuyên gia về lĩnh vực này, nên một chuyên gia có thể chọn những
tập khác. Việc huấn luyện HAN chưa hội tụ hẳn: BPR xác thực còn cải thiện tới \val{han-gain-max}
trong mười epoch cuối, nhưng với các tập tìm được thì nhiều nhất chỉ \val{searched-gain-max}. Do
đó, huấn luyện lâu hơn nhiều khả năng sẽ nới rộng hơn là thu hẹp khoảng cách giữa hai cách. Sản
phẩm chỉ liên kết với các danh mục lá (Mục~\ref{sec:data}). Nếu liên kết mọi danh mục cấp trên, các
ứng viên và lựa chọn của tìm kiếm sẽ khác.

Hướng nghiên cứu tiếp theo gồm: (1) liệu một meta-path danh mục dày đặc có thu hẹp
được khoảng cách không, (2) một phép tìm kiếm được đánh giá bằng chính mô hình gợi ý của nó, như
RMS, (3) các meta-path phía người dùng, (4) truy vấn top-$k$ của \citet{zhu2020topk}, (5) các lược
đồ lớn hơn, các bộ dữ liệu khác và các phép chia theo ngày, nơi tìm kiếm có thể đem lại lợi ích.
